\renewcommand{\chaptertitlename}{Chapter}
\chapter{Star: Update}
\renewcommand{\chaptertitlename}{Capítulo}

\nowplaying{Danza Magia}{Feryquitous}

\lettrine{L}{a} muralla entera parpadea, un latido cian que se apaga por partes, como si algo muy grande estuviera muriendo despacio. Primero los paneles de la izquierda. Después los símbolos tallados que brillaban entre el musgo. Después, justo frente a nosotros, la pantalla de registro.

A mi alrededor, la fila deja de ser fila. Alguien grita que no puede ser, alguien más empuja, alguien llora sin disimularlo. Estamos parados en un saliente de piedra a la altura de un edificio de veinte pisos y el pánico no entiende de alturas.

—¡Por favor, apúrese, señor! —le grita una mujer fea a un desconocido dos puestos más adelante, tirándole de la manga como si eso fuera a arreglar algo. El desconocido no tiene ni idea de qué apurar. Nadie la tiene.

Todos empiezan a enloquecer, cada uno a su manera: unos rezan, unos gritan nombres que no son los suyos, unos simplemente se quedan quietos, mirando el vacío bajo sus pies como si recién ahora lo notaran.

La señora que llevaba el registro —de edad avanzada, manos temblorosas, la única persona en toda la muralla que parecía tener paciencia infinita hasta hace un minuto— se pone de pie y grita algo que no quiero haber escuchado:

—¡No hay sistema! ¡No hay nada que yo pueda hacer! ¡Esta terminal es viejísima, ha administrado todo desde hace generaciones, y nadie, en generaciones, se dignó a cambiarla!

¿Cuántas veces pasó esto antes, en algún otro momento, en algún otro mundo? ¿Cuántas vidas se han perdido a lo largo de la historia no por falta de magia, no por falta de fuerza, sino por seguir confiando en un sistema que ya no servía, solo porque cambiarlo daba pereza, daba miedo, daba trabajo? Pienso en mi propio mundo, el que dejé atrás, el que ni siquiera sé si sigue existiendo. Cuántas veces ahí también preferimos lo roto y conocido antes que lo nuevo e incómodo.

\illustration{images/Chapter-4/socrates-terminal.png}{Sócrates, indiferente al caos, restaura el sistema de la muralla con su propia terminal.}

Entonces aparece él.

No sé de dónde sale —un segundo no está, al siguiente camina entre la gente como si el pánico fuera clima y a él no le afectara el clima. Viejo, descalzo, la túnica cubierta del polvo de mil caminos distintos, con una calma que no debería ser posible a esta altura y con este cielo cayéndose encima. Se detiene un instante frente a todos nosotros y dice, sin levantar la voz, algo que de inmediato sé que voy a recordar el resto de mi vida:

—El secreto para cambiar un sistema roto no es pelear contra él. Es dejar de alimentarlo, y construir, con lo poco que uno tenga a mano, algo que sí funcione.

Y antes de que nadie pueda preguntarle nada, antes de que la sabiduría termine de asentarse en el aire, camina derecho hacia la señora del registro.

—Apártate, mujer, me estorbas.

Ella se aparta. Todos nos apartamos.

De algún pliegue de su túnica saca una terminal pequeña, gastada, con las teclas desgastadas por años de uso real y no de adorno. La conecta al costado de la pantalla muerta con un cable que parece demasiado simple para lo que está a punto de pasar. Escribe. No duda, no corrige, no repite ningún comando. Sus dedos se mueven con la misma economía con la que respira: lo justo, ni un gesto de más.

En menos tiempo del que a mí me tomó subir un solo tramo de esta muralla, los paneles cian vuelven a latir, uno detrás de otro, como un corazón que alguien acaba de reiniciar a la fuerza.

La gente grita, pero esta vez de alivio. La mujer fea abraza al desconocido de la manga. La señora del registro se queda mirando la pantalla encendida con una expresión que no sé si es gratitud o vergüenza.

Sócrates ya no está mirando nada de eso. Guarda su terminal, se da la vuelta, y empieza a caminar otra vez, como si arreglar el sistema que sostenía a cientos de vidas fuera apenas una parada en el camino a otro lado.

León suelta un \textbf{¡HSSSSS!} bajito, y por medio segundo su cara se transforma: cejas altísimas, boca abierta, los ojos brillando con esa mezcla exacta de shock y admiración que normalmente reservo para mí mismo cuando alguien hace algo que yo no sé hacer.

No puedo dejar de mirar la terminal que Sócrates acaba de guardar. Ni un grito, ni un despliegue, ni una sola palabra de más. Solo supo exactamente qué tocar, y lo tocó.

¿Cuántas veces me conformé con un sistema roto solo porque ya estaba acostumbrado a hacer fila frente a él? ¿Cuántas veces elegí la comodidad de lo conocido —aunque no funcionara— antes que la incomodidad de aprender lo que sí funciona?

Tal vez tú también sigas usando, todos los días, un sistema que hace mucho dejó de servirte —una excusa, una costumbre, una versión vieja de ti mismo— solo porque cambiarlo implica aprender algo incómodo. ¿Qué parte tuya necesita, ahora mismo, un simple \emph{reboot}?

No hizo falta magia para salvarnos ahí arriba. Hizo falta alguien que se hubiera tomado el trabajo, en algún momento, de aprender. Eso también se puede aprender. Yo también puedo aprender eso.

El cielo, detrás de todos nosotros, sigue cayendo en pedazos con forma de estrella. Pero por primera vez desde que llegué a este lugar, no tengo tanto miedo de no saber nada. Tengo, en cambio, muchísimas ganas de empezar a saber.
